\chapter{1998:Curtis T. McMullen}\label{ch1998curtistmcmullen:curtis-t-mcmullen-fields-medal-1998-a-comprehensive-report-on-his-work-and-impact}

\begin{center}\rule{0.5\linewidth}{0.5pt}\end{center}

\textbf{Prepared:} July 2026
\textbf{Subject:} Curtis T. McMullen (born April 26, 1958, Berkeley, California)
\textbf{Position:} Professor of Mathematics, Harvard University (since 1997)
\textbf{Doctoral Advisor:} Dennis P. Sullivan (Stony Brook University)
\textbf{Fields Medal Citation:} \emph{``For his contributions to the theory
of 3-manifolds and Kleinian groups.''}

\begin{center}\rule{0.5\linewidth}{0.5pt}\end{center}

\begin{prereqs}
Hyperbolic geometry, Teichmuller theory, complex analysis,
topology of 3-manifolds.
\end{prereqs}

\begin{objectives}
Explain the Bers conjecture and its resolution; describe the
boundary of a hyperbolic 3-manifold; outline the approach to the
ending lamination conjecture.
\end{objectives}

% ============================================================
\section{Problem History}
% ============================================================

\textbf{Hyperbolic 3-manifolds} are quotients
$\mathbb{H}^3/\Gamma$ where $\mathbb{H}^3$ is 3-dimensional
hyperbolic space and $\Gamma\subset\mathrm{PSL}(2,\mathbb{C})$ is
a discrete group of isometries.  The \textbf{Bers conjecture}
(1960) asserted that every hyperbolic 3-manifold with finitely
generated fundamental group has a finite-volume hyperbolic structure.

Thurston's geometrization conjecture (1982) proposed that every
closed orientable 3-manifold can be decomposed into pieces each
admitting one of eight model geometries.  This was partially proved
by Perelman (2003) using Ricci flow, but the hyperbolic structure
of individual pieces required different methods.

A \textbf{Kleinian group} is a discrete subgroup of
$\mathrm{PSL}(2,\mathbb{C})$.  Its \textbf{limit set}
$\Lambda(\Gamma)\subset\widehat{\mathbb{C}}$ is the set of
accumulation points of orbits on the boundary at infinity.  The
\textbf{domain of discontinuity} $\Omega(\Gamma)$ is the
complement.

The \textbf{Bers slice} is the set of Kleinian groups conjugate to
a quasiconformal deformation of a fixed group $\Gamma_0$.  Bers
showed this is parameterized by Teichmuller space, with the
boundary corresponding to degenerate groups.

McMullen's work began in the late 1980s under Sullivan's
mentorship.  Sullivan had established deep connections between
Kleinian groups, quasiconformal mappings, and ergodic theory.
McMullen extended these techniques and developed new ones,
particularly in studying the boundary of hyperbolic 3-manifolds.

The \textbf{Mahler measure} of a polynomial
$P\in\mathbb{Z}[x_1,\dots,x_n]$ is defined as
$M(P) = \exp\bigl(\int_{[0,1]^n}\log|P(e^{2\pi i\theta_1},\dots,
e^{2\pi i\theta_n})|\,d\theta_1\cdots d\theta_n\bigr)$.
McMullen showed that the Mahler measure of certain polynomials is
related to hyperbolic volume, connecting number theory to
3-manifold topology.

% ============================================================
\section{How Was It Solved}
% ============================================================

McMullen's first major contribution was proving \textbf{exponential
volume growth} for a class of hyperbolic 3-manifolds, showing that
the volume of the convex core grows exponentially with the length
of the shortest geodesic, providing quantitative support for the
Bers conjecture.

The \textbf{ending lamination conjecture}, formulated by Thurston,
states that a Kleinian group is determined by its end
invariants---conformal structure at each conformal end and ending
lamination at each geometrically infinite end.  McMullen contributed
by showing the boundary of the convex core is a pleated surface and
developing the theory of \textbf{pleated surfaces} and bending
laminations.

A \textbf{pleated surface} in a hyperbolic 3-manifold is a surface
isometric to a hyperbolic surface except along a lamination where it
bends.  The \textbf{bending measure} on the lamination records how
much the surface bends at each leaf.  McMullen showed that pleated
surfaces provide a canonical geometric representative of surfaces in
hyperbolic 3-manifolds, and that their bending measures control the
geometry of the manifold.

McMullen proved that if $\Lambda(\Gamma)$ is a circle, then
$\Gamma$ is geometrically finite and $\mathbb{H}^3/\Gamma$ has
finite volume.  This is part of understanding when geometric
properties of the limit set imply finiteness conditions.

The key technical innovation was \textbf{renormalization} for
Kleinian groups.  McMullen defined a renormalization map on the
space of Kleinian groups and showed convergence to a fixed point
corresponding to the maximal cusp.

McMullen also developed the theory of \textbf{complex lengths} of
closed geodesics.  A closed geodesic corresponds to a conjugacy
class in $\Gamma$ with complex length $\ell+i\theta$ where $\ell$
is hyperbolic length and $\theta$ is rotation angle.  He showed the
distribution of complex lengths encodes topology and geometry.

The ending lamination conjecture was completed by Minsky (2010),
building on McMullen's techniques.  Together they showed the
mapping class group acts properly discontinuously on bending
laminations.

% ============================================================
\section{Impact of This Discovery}
% ============================================================

McMullen's work on hyperbolic 3-manifolds has been central to
3-manifold topology.  The proof that geometrically infinite
Kleinian groups have exotic limit sets influenced the
classification of 3-manifolds.

The theory of pleated surfaces and bending laminations provides
the geometric framework for understanding hyperbolic 3-manifold
boundaries, applied to essential surfaces, the Smith conjecture,
and knot classification.

The ending lamination conjecture provides a complete
parameterization of Kleinian groups, establishing a dictionary
between 3-manifold topology and hyperbolic geometry.

McMullen's work on rigidity influenced the study of arithmetic
Kleinian groups.  The application of holomorphic dynamics to
Kleinian groups created fruitful interaction between previously
separate areas.

The \textbf{A-polynomial} of a knot, introduced by Cooper, Culler,
Gillet, Long, and Shalen, encodes the character variety of a knot
complement.  McMullen's work on complex lengths provides geometric
information about this character variety.

McMullen's influence extends beyond hyperbolic geometry.  His work
on \textbf{Riemann surfaces and their moduli} has been fundamental,
including results on the Teichmuller geodesic flow, the Lattes
maps, and the classification of post-critically finite rational
maps.  His monograph ``Complex Dynamics and Renormalization'' is a
standard reference in holomorphic dynamics.

% ============================================================
\section{Modern Views and Improvements}
% ============================================================

The ending lamination conjecture, proved by Minsky, established the
Bers parameterization.  McMullen's contributions remain fundamental.

The volume conjecture for knots connects the colored Jones
polynomial to hyperbolic volume.  McMullen's work on volume
distributions has been relevant.

Random Kleinian groups, initiated by McMullen, has become active.
Hausdorff dimension of limit sets connects probability to
hyperbolic geometry.

The virtual Haken conjecture, proved by Agol, builds on hyperbolic
geometry that McMullen and others developed.  The Agol--Wilton
proof provides complete resolution of Thurston's geometrization.

Recent work by Brock and Canary shows that the \textbf{boundary of
the convex core} of a Kleinian group has Hausdorff dimension
related to the Weil--Petersson metric on Teichmuller space,
extending McMullen's geometric analysis.

The \textbf{Cannon--Thurston map} for Kleinian groups, which
extends the natural map from the universal cover of a surface to
the boundary at infinity to a continuous map of the sphere, has been
studied extensively using McMullen's techniques.

The \textbf{Minsky machine} algorithm computes the hyperbolic
3-manifold from its end invariants in polynomial time, effectively
making the ending lamination theorem constructive.

% ============================================================
\begin{center}\rule{0.5\linewidth}{0.5pt}\end{center}

\section*{Summary}
\addcontentsline{toc}{section}{Summary}

McMullen made fundamental contributions to hyperbolic 3-manifolds
and Kleinian groups, including renormalization techniques, pleated
surface theory, and essential contributions to the ending lamination
conjecture.  His work connected holomorphic dynamics to 3-manifold
topology.

% ============================================================
